\documentclass[11pt]{article}
\usepackage{amsmath,amssymb,amsthm}
\usepackage[margin=2.5cm]{geometry}

\newtheorem{theorem}{Theorem}
\newtheorem{lemma}[theorem]{Lemma}
\theoremstyle{remark}
\newtheorem*{remark}{Remark}

\newcommand{\R}{\mathbb{R}}

\title{Sector decomposition for an anisotropic oscillatory integral}
\author{}
\date{}

\begin{document}
\maketitle

\section{Setting and overview}

\subsection{The integral}

Let $\ell_1,\dots,\ell_d\ge0$ and, for $s>0$, let
$A_s=\operatorname{diag}(s^{\ell_1},\dots,s^{\ell_d})$. For $t\ge1$ and $\xi\in\R^d$
we define the phase
\[
\psi(t,\xi)=\int_1^t|A_s\xi|\,ds .
\]
We study the oscillatory integral
\[
I(t,x)=\int_{\R^d}e^{i(x\cdot\xi-\psi(t,\xi))}\,|A_t\xi|^{-1}\,\chi(\xi)\,d\xi ,
\]
where $\chi$ is a smooth cutoff supported in the annulus $\{\tfrac12\le|\xi|\le2\}$.
We want to bound $I(t,x)$ for large $t$, uniformly in $x$.

\subsection{The difficulty and the plan}

The matrix $A_t$ stretches each frequency coordinate by a different power of $t$.
So there is no single large parameter in front of the phase. We handle this in two
moves.

First we split frequency space into $d$ sectors, according to which weighted
coordinate $t^{\ell_j}|\xi_j|$ is the largest. On sector $k$ we use the dominant
coordinate $\xi_k$ and the $d-1$ ratios $\eta_j=t^{\ell_j}\xi_j/(t^{\ell_k}\xi_k)$ as
new variables. In these variables the sector integral becomes an outer integral in
$\xi_k$ of an inner oscillatory integral in $\eta$ with a single large parameter
$\lambda=t^{\ell_k+1}|\xi_k|$ (Section~\ref{sec:reduction}).

Second we bound the inner integral by stationary phase. The phase is strictly
concave in $\eta$, and its curvature in each direction is at least of order $1$; in
borderline directions it reaches order $\log t$ when the leading coordinate takes
over early enough (Section~\ref{sec:hessian}). Its
higher derivatives are not bounded uniformly in $t$, but they are once we weight
them correctly (Section~\ref{sec:higher}). The stationary phase theorem needs its
constants to be bounded uniformly in $t$. At unit scale in $\eta$ they are not; at
the natural scale of each direction, $(\lambda L_j)^{-1/2}$ with $L_j=1$ or
$L_j=\log t$, they are. Section~\ref{sec:spt} checks this, and
Section~\ref{sec:proof} assembles the main result, Theorem~\ref{thm:main}. The outer
integral in $\xi_k$ is not treated here.

\subsection{Conventions}

In all estimates $t\ge2$. The index $k$ is fixed, $n=d-1$, and $i,j,m$ run over
$\{1,\dots,d\}\setminus\{k\}$. Multi-indices $\alpha=(\alpha_j)_{j\ne k}$ have $d-1$
entries, and $\delta_i$ is the multi-index with a $1$ in position $i$ and $0$
elsewhere. We write $C$, $c$ for positive constants that depend only on $d$ and
$\ell_1,\dots,\ell_d$ and may change from line to line. $C_\alpha$ may also depend
on $\alpha$.

\section{Reduction to one sector}\label{sec:reduction}

\subsection{Sectors}

Sector $k$ is the set of $\xi$ for which $t^{\ell_k}|\xi_k|$ is the largest of the
weights $t^{\ell_1}|\xi_1|,\dots,t^{\ell_d}|\xi_d|$. Since $|A_t\xi|^2$ is the sum
of the $d$ squared weights, on sector $k$
\[
t^{\ell_k}|\xi_k|\ \ge\ \frac{|A_t\xi|}{\sqrt d}.
\]
Every $\xi$ lies in some sector, so $I$ is a sum of $d$ sector integrals. We fix $k$.

\subsection{New variables}

For $j\ne k$ we introduce the ratios
\[
\eta_j=\frac{t^{\ell_j}\xi_j}{t^{\ell_k}\xi_k},
\]
which compare the weight of coordinate $j$ with the dominant weight. On sector $k$
they satisfy $|\eta_j|\le1$. We write $\eta=(\eta_j)_{j\ne k}\in\R^{d-1}$ and use
$(\xi_k,\eta)$ as coordinates on the sector.

\subsection{The integrand in the new variables}

We rewrite each factor of the integrand in terms of $\xi_k$ and $\eta$.

\paragraph{The amplitude.} Since
\[
|A_t\xi|^2=t^{2\ell_k}\xi_k^2+\sum_{j\ne k}t^{2\ell_j}\xi_j^2
=t^{2\ell_k}\xi_k^2\Big(1+\sum_{j\ne k}\eta_j^2\Big),
\]
we have
\[
|A_t\xi|=t^{\ell_k}|\xi_k|\,(1+|\eta|^2)^{1/2},
\qquad
|A_t\xi|^{-1}=t^{-\ell_k}|\xi_k|^{-1}(1+|\eta|^2)^{-1/2}.
\]

\paragraph{The phase $\psi$.} In the same way, for $0<s\le t$,
\[
|A_s\xi|=s^{\ell_k}|\xi_k|\Big(1+\sum_{j\ne k}\frac{s^{2\ell_j}\xi_j^2}{s^{2\ell_k}\xi_k^2}\Big)^{1/2}.
\]
We measure time in units of $t$: substitute $s=\sigma t$, so $ds=t\,d\sigma$ and
$s^{2(\ell_j-\ell_k)}\xi_j^2/\xi_k^2=\sigma^{2(\ell_j-\ell_k)}\eta_j^2$. This gives
\[
\psi(t,\xi)=t^{\ell_k+1}|\xi_k|\,G(t,\eta),
\qquad
G(t,\eta)=\int_{1/t}^1\sigma^{\ell_k}\Big(1+\sum_{j\ne k}\sigma^{2(\ell_j-\ell_k)}\eta_j^2\Big)^{1/2}d\sigma .
\]
So the phase is a large factor $t^{\ell_k+1}|\xi_k|$ times a function $G$ of
$\eta$, whose only dependence on $t$ is through the lower limit $1/t$.

\paragraph{The linear term.} We aim at a bound uniform in $x$, so we may replace $x$
by $A_tx$ without changing $\sup_x|I(t,x)|$. The linear part of the phase becomes
\[
A_tx\cdot\xi=x_kt^{\ell_k}\xi_k+\sum_{j\ne k}t^{\ell_j}x_j\xi_j
=x_kt^{\ell_k}\xi_k+t^{\ell_k}\xi_k\sum_{j\ne k}x_j\eta_j .
\]

\paragraph{The measure.} For fixed $\xi_k$, the change of variables
$\xi_j=t^{\ell_k-\ell_j}\xi_k\eta_j$, $j\ne k$, gives
\[
d\xi=|\xi_k|^{d-1}\prod_{j\ne k}t^{\ell_k-\ell_j}\;d\xi_k\,d\eta
=t^{\,d\ell_k-\sum_{j=1}^d\ell_j}\,|\xi_k|^{d-1}\,d\xi_k\,d\eta .
\]

\subsection{The sector integral}\label{sec:sector-integral}

Denote by $I_k$ the part of $I$ coming from sector $k$. Collecting the formulas
above,
\[
I_k=t^{\,d\ell_k-\sum_j\ell_j-\ell_k}\int_{\R}e^{ix_kt^{\ell_k}\xi_k}\,|\xi_k|^{d-2}\,
J(\xi_k)\,d\xi_k ,
\qquad
J(\xi_k)=\int_{[-1,1]^{d-1}}e^{i\lambda\Phi(\eta)}\,b(\xi_k,\eta)\,d\eta ,
\]
where the large parameter, the phase and the amplitude are
\[
\lambda=t^{\ell_k+1}|\xi_k| ,
\qquad
\Phi(\eta)=\frac{\xi_k}{|\xi_k|}\,\frac1t\sum_{j\ne k}x_j\eta_j-G(t,\eta),
\]
\[
b(\xi_k,\eta)=(1+|\eta|^2)^{-1/2}\,
\chi\Big(\xi_ke_k+\sum_{j\ne k}t^{\ell_k-\ell_j}\xi_k\eta_j\,e_j\Big).
\]
The rest of these notes bounds the inner integral $J$. Two remarks on its
ingredients.

\paragraph{The phase.} $\Phi$ depends on $x$ and on the sign of $\xi_k$ only
through its linear part. All derivatives of order two and higher are those of $-G$
and do not see $x$ at all.

\paragraph{The amplitude and the size of $\lambda$.} The second factor of $b$ is the
cutoff $\chi$ in the new variables. Since
$|\xi|=|\xi_k|\big(1+\sum_{j\ne k}t^{2(\ell_k-\ell_j)}\eta_j^2\big)^{1/2}$, the
amplitude $b$ is supported where
\[
\frac12\le|\xi_k|\Big(1+\sum_{j\ne k}t^{2(\ell_k-\ell_j)}\eta_j^2\Big)^{1/2}\le2 .
\]
With $\ell_{\min}=\min_j\ell_j$ and $|\eta_j|\le1$, this gives
\[
\frac{t^{\ell_{\min}-\ell_k}}{2\sqrt d}\ \le\ |\xi_k|\ \le\ 2
\qquad\text{on the support of } b,
\]
and therefore $\lambda\ge t^{1+\ell_{\min}}/(2\sqrt d)$. So $\lambda$ is indeed
large. Note, however, that a derivative $\partial_{\eta_j}$ falling on the cutoff
produces a factor $t^{\ell_k-\ell_j}\xi_k$, which is large when $\ell_j<\ell_k$.
This will be the main obstacle in Section~\ref{sec:spt}.

\subsection{Types of directions}\label{sec:types}

How each direction $j\ne k$ is treated depends on the number
\[
\kappa_j=2\ell_j-\ell_k+1 .
\]
We call direction $j$
\begin{itemize}
	\item \emph{good} if $\ell_j\ge\ell_k$, and \emph{bad} if $\ell_j<\ell_k$;
	\item \emph{oscillatory} if $\kappa_j>0$, \emph{borderline} if $\kappa_j=0$, and
	\emph{flat} if $\kappa_j<0$.
\end{itemize}
We write $\mathcal O$, $\mathcal Z$ and $\mathcal F$ for the sets of oscillatory,
borderline and flat directions, $\mathcal T=\mathcal Z\cup\mathcal F$ for the
non-oscillatory ones, and $\mathcal B$ for the bad oscillatory ones. Every good
direction is oscillatory, since then $\kappa_j\ge\ell_k+1>0$. So borderline and flat
directions are bad. If there is a bad direction we put
\[
\theta=\tfrac12\min_{m\text{ bad}}(\ell_k-\ell_m)>0 .
\]
In short: oscillatory directions gain from oscillation, flat directions gain from
the small size of their support, and borderline directions gain from oscillation
with an extra factor $(\log t)^{-1/2}$ when the curvature has time to build up.

\section{The Hessian of the phase}\label{sec:hessian}

Set
\[
H(\sigma,\eta)=\Big(1+\sum_{j\ne k}\sigma^{2(\ell_j-\ell_k)}\eta_j^2\Big)^{1/2},
\qquad\text{so that}\qquad
G(t,\eta)=\int_{1/t}^1\sigma^{\ell_k}\,H(\sigma,\eta)\,d\sigma .
\]
The goal of this section is to show that $\Phi$ is strictly concave in $\eta$, and
to measure its curvature in each direction.

\paragraph{First derivatives.} Since
$\partial_{\eta_j}H=\sigma^{2(\ell_j-\ell_k)}\eta_j\,H^{-1}$,
\[
\partial_{\eta_j}\Phi(\eta)=\frac{\xi_k}{|\xi_k|}\,\frac{x_j}{t}
-\eta_j\int_{1/t}^1\sigma^{2\ell_j-\ell_k}\,H^{-1}\,d\sigma .
\]

\paragraph{Second derivatives.} The linear part of $\Phi$ has zero Hessian, so
$\nabla_\eta^2\Phi=-\nabla_\eta^2G$. Differentiating once more,
\[
\partial_{\eta_i}\partial_{\eta_j}G(t,\eta)
=\int_{1/t}^1\sigma^{\ell_k}\,H^{-3}
\Big(\delta_{ij}\,\sigma^{2(\ell_j-\ell_k)}\,H^2
-\sigma^{2(\ell_i+\ell_j-2\ell_k)}\eta_i\eta_j\Big)\,d\sigma .
\]

\paragraph{The quadratic form.} Let
$D(\sigma)=\operatorname{diag}\big(\sigma^{2(\ell_j-\ell_k)}\big)_{j\ne k}$. In matrix
form the second derivatives read
\[
-\nabla_\eta^2\Phi
=\int_{1/t}^1\sigma^{\ell_k}H^{-3}\Big(H^2D-(D\eta)(D\eta)^{T}\Big)\,d\sigma .
\]
Fix $v\in\R^{d-1}$ and put $z=D^{1/2}\eta$, $u=D^{1/2}v$. Then
$\langle Dv,v\rangle=|u|^2$, $\langle D\eta,v\rangle=\langle z,u\rangle$, and
$|z|^2=\langle D\eta,\eta\rangle=H^2-1$. Hence
\[
-\langle\nabla_\eta^2\Phi\,v,v\rangle
=\int_{1/t}^1\sigma^{\ell_k}H^{-3}\Big(H^2|u|^2-\langle z,u\rangle^2\Big)\,d\sigma .
\]
By Cauchy--Schwarz, $\langle z,u\rangle^2\le|z|^2|u|^2=(H^2-1)|u|^2$, so the
bracket lies between $|u|^2$ and $H^2|u|^2$. Writing
\[
S(\sigma,v)=|u|^2=\sum_{j\ne k}\sigma^{2(\ell_j-\ell_k)}v_j^2 ,
\]
we arrive at the two-sided bound
\[
\int_{1/t}^1\sigma^{\ell_k}H^{-3}\,S\,d\sigma
\;\le\;-\langle\nabla_\eta^2\Phi\,v,v\rangle\;\le\;
\int_{1/t}^1\sigma^{\ell_k}H^{-1}\,S\,d\sigma .
\]

\paragraph{Strict concavity.} For $v\ne0$ the left side is positive. So
$-\nabla_\eta^2\Phi$ is positive definite and $\Phi$ is strictly concave in $\eta$,
for every $t$, $x$ and $\xi_k$.

\paragraph{The curvature in each direction.} The lower bound is diagonal. Since
$\sigma^{\ell_k}\sigma^{2(\ell_j-\ell_k)}=\sigma^{2\ell_j-\ell_k}=\sigma^{\kappa_j-1}$,
it reads
\[
-\langle\nabla_\eta^2\Phi\,v,v\rangle\ \ge\ \sum_{j\ne k}\Lambda_j(\eta)\,v_j^2,
\qquad
\Lambda_j(\eta)=\int_{1/t}^1\sigma^{\kappa_j-1}\,H(\sigma,\eta)^{-3}\,d\sigma .
\]
On a dyadic time interval $\sigma\in[2^{-i-1},2^{-i}]$ where $H$ stays bounded, the
integrand of $\Lambda_j$ contributes about $2^{-i\kappa_j}$. So the curvature of an
oscillatory direction is built at late times, that of a borderline direction equally
at every time scale, and that of a flat direction mostly at early times. To bound
$\Lambda_j$ from below we need to know when $H$ is bounded.

\paragraph{The switch time.} For a bad direction $m$, the term
$\sigma^{2(\ell_m-\ell_k)}\eta_m^2$ of $H^2$ is at most $1$ exactly when
$\sigma\ge|\eta_m|^{1/(\ell_k-\ell_m)}$. We call
\[
\sigma_*(\eta)=\max_{m\text{ bad}}|\eta_m|^{1/(\ell_k-\ell_m)}
\]
the switch time, with $\sigma_*=0$ if there is no bad direction. In the original
variables, $\sigma^{\ell_m-\ell_k}|\eta_m|=s^{\ell_m}|\xi_m|/(s^{\ell_k}|\xi_k|)$ with
$s=\sigma t$. So $s_*=\sigma_*t$ is the time at which the weight $s^{\ell_k}|\xi_k|$
overtakes every bad weight $s^{\ell_m}|\xi_m|$.

\begin{lemma}[Curvature]\label{lem:curv}
	Let $R\ge1$, $t\ge2$ and $|\eta|\le R$. Then for all $v\in\R^{d-1}$,
	\[
	-\langle\nabla_\eta^2\Phi(\eta)\,v,v\rangle\ \ge\ c_R\sum_{j\ne k}\Gamma_j(\eta)\,v_j^2,
	\qquad
	\Gamma_j(\eta)=\int_{\sigma_0}^1\sigma^{\kappa_j-1}\,d\sigma ,
	\]
	where $\sigma_0=\max\big(1/t,\min(\sigma_*(\eta),\tfrac12)\big)$ and $c_R>0$
	depends only on $R$, $d$ and the $\ell_j$. In particular:
	\begin{enumerate}
		\item[(i)] $\Gamma_j\ge c$ for every $j$, so $-\nabla_\eta^2\Phi\ge c_R\,I$ on
		$|\eta|\le R$;
		\item[(ii)] if $j$ is borderline and $\sigma_*(\eta)\le t^{-1/2}$, then
		$\Gamma_j(\eta)\ge\frac12\log t$.
	\end{enumerate}
\end{lemma}

\begin{proof}
	We show that $H\le C_R$ on $[\sigma_0,1]$. A good $m$ contributes
	$\sigma^{2(\ell_m-\ell_k)}\eta_m^2\le R^2$. A bad $m$ contributes at most $1$ when
	$\sigma\ge\sigma_*$, and at most $4^{\ell_k-\ell_m}R^2$ when $\sigma\ge\frac12$; one
	of the two holds on $[\sigma_0,1]$. So $H^{-3}\ge c_R$ there, and since the
	integrand of $\Lambda_j$ is nonnegative, $\Lambda_j\ge c_R\Gamma_j$.
	For (i), $\sigma_0\le\frac12$, so $\Gamma_j\ge\int_{1/2}^1\sigma^{\kappa_j-1}d\sigma$,
	a positive constant. For (ii), $\kappa_j=0$ gives $\Gamma_j=\log(1/\sigma_0)$, and
	$\sigma_0\le\max(1/t,t^{-1/2})=t^{-1/2}$.
\end{proof}

Part (ii) is where $\log t$ enters. A borderline direction collects the same amount
$\log2$ of curvature from every dyadic time scale after the switch. If the switch
comes before $t^{1/2}$ in physical time, about $\frac12\log_2t$ scales contribute.
For flat directions $\Gamma_j$ is even larger, but there the amplitude is too rough
at the stationary phase scale, and this curvature is not used.

\paragraph{An upper bound.} Conversely, the upper half of the two-sided bound and
$H\ge1$ give
\[
\partial_{\eta_j}^2G\ \le\ \bar\Lambda_j=\int_{1/t}^1\sigma^{\kappa_j-1}\,d\sigma ,
\]
with $\bar\Lambda_j\le1/\kappa_j$ for oscillatory $j$ and $\bar\Lambda_j=\log t$ for
borderline $j$. So after an early switch, the curvature of a borderline direction is
exactly of order $\log t$.

\section{Higher derivatives of the phase}\label{sec:higher}

Let $|\alpha|\ge2$. Two derivatives kill the linear part of $\Phi$, so
$\partial_\eta^\alpha\Phi=-\partial_\eta^\alpha G$, and it suffices to bound
$\partial_\eta^\alpha G$. We first give a direct bound. Then we explain why it is
not uniform in $t$, and we prove a weighted bound that is.

\subsection{A direct bound}

\paragraph{Reduction to one function.} For fixed $\sigma$ put
$y_j=\sigma^{\ell_j-\ell_k}\eta_j$, $j\ne k$, and $\langle y\rangle=(1+|y|^2)^{1/2}$.
Then $H(\sigma,\eta)=\langle y\rangle$. Since each $y_j$ is a multiple of $\eta_j$
alone, the chain rule gives
\[
\partial_\eta^\alpha H
=\prod_{j\ne k}\sigma^{\alpha_j(\ell_j-\ell_k)}\;\big(\partial_y^\alpha\langle y\rangle\big).
\]

\paragraph{Derivatives of $\langle y\rangle$.} We claim that for $|\alpha|\ge1$,
$\partial_y^\alpha\langle y\rangle$ is a finite linear combination of terms
\[
y^\beta\,\langle y\rangle^{1-2m}\qquad\text{with}\qquad m\ge1,\quad2m-|\beta|=|\alpha| .
\]
For $|\alpha|=1$ this holds, since $\partial_{y_i}\langle y\rangle=y_i\langle y\rangle^{-1}$.
Differentiating $y^\beta\langle y\rangle^{1-2m}$ in $y_i$ gives two terms:
$\beta_i\,y^{\beta-\delta_i}\langle y\rangle^{1-2m}$, which vanishes if $\beta_i=0$,
and $(1-2m)\,y^{\beta+\delta_i}\langle y\rangle^{1-2(m+1)}$. In both terms
$2m-|\beta|$ increases by one, which proves the claim. Since
$|y^\beta|\le\langle y\rangle^{|\beta|}$, each term is bounded by
$\langle y\rangle^{|\beta|+1-2m}=\langle y\rangle^{1-|\alpha|}$, so
\[
\big|\partial_y^\alpha\langle y\rangle\big|\le C_\alpha\,\langle y\rangle^{1-|\alpha|} .
\]
In words, $\langle y\rangle$ behaves like a function homogeneous of degree one.

\paragraph{The bound.} Differentiating under the integral sign and using
$\langle y\rangle=H\ge1$, we obtain for $|\alpha|\ge1$
\[
\big|\partial_\eta^\alpha G(t,\eta)\big|
\le C_\alpha\int_{1/t}^1\sigma^{\,\ell_k+\sum_{j\ne k}\alpha_j(\ell_j-\ell_k)}\,
H(\sigma,\eta)^{1-|\alpha|}\,d\sigma
\le C_\alpha\int_{1/t}^1\sigma^{\,\ell_k+\sum_{j\ne k}\alpha_j(\ell_j-\ell_k)}\,d\sigma .
\]
The last integral stays bounded as $t\to\infty$ only if the exponent of $\sigma$ is
larger than $-1$. When some $\ell_j<\ell_k$, this fails for large $\alpha_j$. We
look at this more closely.

\subsection{Bad directions and weights}

Recall from Section~\ref{sec:types} that direction $j$ is bad if $\ell_j<\ell_k$.

\paragraph{Why the direct bound fails.} High derivatives in a bad direction blow up
at $\eta_j=0$. Indeed, set all $\eta_m$ with $m\ne j$ to zero and expand
$H=(1+\sigma^{2(\ell_j-\ell_k)}\eta_j^2)^{1/2}$ in powers of $\eta_j$. The
coefficient of $\eta_j^{2N}$ is $\binom{1/2}{N}\sigma^{2N(\ell_j-\ell_k)}$, so
\[
\partial_{\eta_j}^{2N}G(t,0)=(2N)!\binom{1/2}{N}\int_{1/t}^1\sigma^{\,\ell_k+2N(\ell_j-\ell_k)}\,d\sigma .
\]
Since $\ell_j<\ell_k$, the exponent is below $-1$ for $N$ large, and the integral
tends to infinity as $t\to\infty$. The reason is the switch of Section~\ref{sec:hessian}:
near $\eta_j=0$ the coordinate $j$ dominates $H$ at early times $\sigma\sim1/t$,
where $\sigma^{\ell_j-\ell_k}$ is huge. Away from $\eta_j=0$ this does not happen.

\paragraph{The weights.} We measure derivatives in units of the distance to these
hyperplanes:
\[
v_j(\eta)=1\ \ (j\text{ good}),\qquad v_j(\eta)=|\eta_j|\ \ (j\text{ bad}).
\]

\subsection{The weighted bound}

\paragraph{Derivatives of $H$, weighted.} Write $\mu_m=\sigma^{\ell_m-\ell_k}$, so
that $y_m=\mu_m\eta_m$. The chain rule and the bound for $\langle y\rangle$ give
\[
\big|\partial_\eta^\alpha H\big|\le C_\alpha\,\mu^\alpha\,H^{1-|\alpha|}
=C_\alpha\,H\prod_{m\ne k}\Big(\frac{\mu_m}{H}\Big)^{\alpha_m},
\qquad\mu^\alpha=\prod_{m\ne k}\mu_m^{\alpha_m}.
\]
Two simple facts control each factor $\mu_m/H$:
\begin{itemize}
	\item if $m$ is good, then $\mu_m\le1\le H$, so $\mu_m/H\le1$;
	\item if $m$ is bad, then $H\ge\mu_m|\eta_m|$, so $\mu_m/H\le|\eta_m|^{-1}$.
\end{itemize}
In both cases $\mu_m/H\le v_m^{-1}$.

\begin{lemma}[Weighted derivatives]\label{lem:weighted}
	Let $|\alpha|\ge2$ with $\alpha_m=0$ for $m\notin\mathcal O\cup\mathcal Z$, and
	choose $i,j\in\mathcal O\cup\mathcal Z$ (possibly $i=j$) such that
	$\alpha'=\alpha-\delta_i-\delta_j$ has nonnegative entries. Then for $t\ge2$,
	\[
	\big|\partial_\eta^\alpha G(t,\eta)\big|\le C_\alpha\,\big(\bar\Lambda_i\bar\Lambda_j\big)^{1/2}
	\prod_{m\ne k}v_m(\eta)^{-\alpha'_m},
	\]
	where $\bar\Lambda_m=\int_{1/t}^1\sigma^{\kappa_m-1}d\sigma$ is at most $1/\kappa_m$
	for $m\in\mathcal O$ and equals $\log t$ for $m\in\mathcal Z$.
\end{lemma}

In words: two derivatives in oscillatory directions are free, a borderline one among
them costs at most $(\log t)^{1/2}$, and each further derivative in a bad direction
costs a factor $1/|\eta_m|$. Further derivatives in good directions cost nothing.

\begin{proof}
	Since $|\alpha|=|\alpha'|+2$, we can split the bound for $\partial_\eta^\alpha H$ as
	\[
	\mu^\alpha H^{1-|\alpha|}=\frac{\mu_i\mu_j}{H}\,\prod_{m\ne k}\Big(\frac{\mu_m}{H}\Big)^{\alpha'_m}
	\le\frac{\mu_i\mu_j}{H}\,\prod_{m\ne k}v_m^{-\alpha'_m}.
	\]
	The product on the right does not depend on $\sigma$. It remains to integrate the
	first factor against $\sigma^{\ell_k}d\sigma$. Since $\sigma^{\ell_k}\mu_m^2
	=\sigma^{\kappa_m-1}$, the bound $H\ge1$ and Cauchy--Schwarz give
	\[
	\int_{1/t}^1\sigma^{\ell_k}\,\frac{\mu_i\mu_j}{H}\,d\sigma
	\le\Big(\int_{1/t}^1\sigma^{\kappa_i-1}d\sigma\Big)^{1/2}
	\Big(\int_{1/t}^1\sigma^{\kappa_j-1}d\sigma\Big)^{1/2}
	=\big(\bar\Lambda_i\bar\Lambda_j\big)^{1/2}.
	\]
	Finally $\bar\Lambda_m=(1-t^{-\kappa_m})/\kappa_m\le1/\kappa_m$ if $\kappa_m>0$,
	and $\bar\Lambda_m=\log t$ if $\kappa_m=0$.
\end{proof}

\subsection{What the weights suggest}

Lemma~\ref{lem:weighted} is made for the scale of stationary phase. At the scale
$\lambda^{-1/2}$, a further derivative in a bad direction $m$ costs
$\lambda^{-1/2}/|\eta_m|$, which is at most $1$ as soon as
$|\eta_m|\ge\lambda^{-1/2}$. So the weighted bound controls all derivatives except
in thin strips around the hyperplanes $\eta_m=0$. A strip has width about
$\lambda^{-1/2}$, exactly the gain oscillation would give, so it can be integrated
trivially. Section~\ref{sec:spt} carries this out.

\section{The stationary phase theorem at the zoom scale}\label{sec:spt}

\subsection{Smooth sectors}

The stationary phase estimate below needs a smooth amplitude, and the indicator of
a sector is not smooth. So we replace the sharp sectors by a smooth partition. Fix
smooth $\rho_1,\dots,\rho_d$ on $\mathbb S^{d-1}$ with $\sum_k\rho_k=1$ and
$\rho_k(p)=0$ unless $|p_k|\ge\frac1{2\sqrt d}$. Such a partition exists
since $\max_k|p_k|\ge1/\sqrt d$ on $\mathbb S^{d-1}$. We insert
$\rho_k(A_t\xi/|A_t\xi|)$ into $I_k$ in place of the indicator of sector $k$. The
$\eta$ integral then runs over $\R^{d-1}$, with amplitude
\[
b(\xi_k,\eta)=(1+|\eta|^2)^{-1/2}\,\tilde\chi(\xi_k,\eta)\,
\rho_k\Big(\frac{\xi_k}{|\xi_k|}\,\frac{e_k+\sum_{j\ne k}\eta_je_j}{(1+|\eta|^2)^{1/2}}\Big),
\]
where
$\tilde\chi(\xi_k,\eta)=\chi\big(\xi_ke_k+\sum_{j\ne k}t^{\ell_k-\ell_j}\xi_k\eta_j\,e_j\big)$
is the cutoff in the new variables. The $k$th component of the argument of $\rho_k$
is $\pm(1+|\eta|^2)^{-1/2}$, so $b(\xi_k,\cdot)$ is smooth and supported in
$|\eta|\le2\sqrt d$. The bounds of Section~\ref{sec:sector-integral} on the support
remain true with other constants: $|\xi_k|\le2$ and $\lambda\ge c\,t^{1+\ell_{\min}}$.
Moreover $|\xi_m|\le2$ for every $m$ on the support of $\tilde\chi$, that is,
\begin{equation}\label{eq:support}
|\eta_m|\le\frac{2\,t^{\ell_m-\ell_k}}{|\xi_k|}\qquad(m\ne k).
\end{equation}

\subsection{The theorem}

For $b\in C_0^\infty(\R^n)$ with support $K$, let $K_1$ be the set of points at sup
distance at most $1$ from $K$, and set
\[
M_m=\sum_{2\le|\alpha|\le m}\ \sup_{K_1}|\partial^\alpha\Phi|,\qquad
N_m=\sum_{|\alpha|\le m}\ \sup_K|\partial^\alpha b|,\qquad
a_0=\inf_{K_1}\big|\det\nabla^2\Phi\big| .
\]

\begin{theorem}[Alazard, Burq, Zuily, Theorem 3]\label{thm:ABZ}
	Let $\Phi$ be real and smooth near $K_1$, with $a_0>0$ and $\nabla\Phi$ injective
	on $K_1$. Then for all $\lambda\ge1$,
	\[
	\Big|\int_{\R^n}e^{i\lambda\Phi(\eta)}\,b(\eta)\,d\eta\Big|
	\le\frac{C_n}{a_0}\,\big(1+M_{n+2}\big)^{n/2}\,N_{n+1}\,\lambda^{-n/2},
	\]
	where $C_n$ depends only on $n$.
\end{theorem}

Two features of this bound matter for us. It involves no volume of $K$, so a large
support costs nothing. And it sees $\Phi$ only through derivatives of order at
least two, so the linear part of $\Phi$, which carries $x$, plays no role.

\subsection{Why the unit scale is not enough}

Applied directly in $\eta$, Theorem~\ref{thm:ABZ} has its hypotheses satisfied for
every fixed $t$: $a_0\ge c$ by Lemma~\ref{lem:curv}, and $\nabla\Phi$ is injective by
concavity. But its constant is not uniform in $t$. Each $\partial_{\eta_j}$ falling
on $\tilde\chi$ produces the factor $t^{\ell_k-\ell_j}\xi_k$, so $N_d$ can be as large
as $t^{d(\ell_k-\ell_{\min})}$. By Section~\ref{sec:higher}, the phase derivatives in
$M_{d+1}$ grow with $t$ in bad directions. Borderline directions add a milder
problem: Lemma~\ref{lem:curv} can give $a_0$ a factor $\log t$ for each of them, but
their second derivatives in $M$ are then of size $\log t$ as well, and the factor
$(1+M)^{n/2}$ can cancel the gain.

All of this comes from measuring derivatives at unit scale in $\eta$, while the
integral lives at scale $\lambda^{-1/2}$, and at scale $(\lambda\log t)^{-1/2}$ in
borderline directions. We now check the hypotheses at that scale.

\subsection{The zoom scale}\label{sec:zoom}

\paragraph{Early and late switch.} If there is a bad direction, we say that the
switch is \emph{early} if $|\xi_k|\ge4t^{-\theta}$, and \emph{late} otherwise. By
\eqref{eq:support} and the switch time of Section~\ref{sec:hessian}, an early switch
means that the weight $s^{\ell_k}|\xi_k|$ overtakes every bad weight before the time
$s=t^{1/2}$; this is made precise in part (a) of the proof below. If there is no bad
direction, there is no borderline direction either, and the distinction plays no
role.

\paragraph{Active candidates and scales.} The directions used for stationary phase
are
\[
\mathcal A=\mathcal O\cup\mathcal Z\ \ \text{(early switch)},\qquad
\mathcal A=\mathcal O\ \ \text{(late switch)}.
\]
Flat directions, and borderline directions after a late switch, are never used for
stationary phase. The curvature scale of direction $j$ is
\[
L_j=\log t\ \ \text{if $j\in\mathcal Z$ and the switch is early},\qquad
L_j=1\ \ \text{otherwise}.
\]
We take $t_0\ge e$, so that $L_j\ge1$.

\paragraph{The zoom.} Put $\eta_j=(\lambda L_j)^{-1/2}\omega_j$ for every $j\ne k$.
Then
\[
J=\lambda^{-n/2}\prod_{j\ne k}L_j^{-1/2}\int_{\R^n}e^{i\tilde\Phi(\omega)}\,\tilde b(\omega)\,d\omega ,
\qquad
\tilde\Phi(\omega)=\lambda\,\Phi(\eta),
\quad
\tilde b(\omega)=b(\xi_k,\eta),
\]
and by the chain rule
\[
\partial_\omega^\alpha\tilde\Phi=\lambda^{1-|\alpha|/2}\prod_{m\ne k}L_m^{-\alpha_m/2}\,
\partial_\eta^\alpha\Phi .
\]
For $|\alpha|=2$ this reads
$\partial_{\omega_i}\partial_{\omega_j}\tilde\Phi=(L_iL_j)^{-1/2}\partial_{\eta_i}\partial_{\eta_j}\Phi$:
the Hessian is normalized by the curvature scales. The decay $\lambda^{-n/2}$ and the
gain $\prod_jL_j^{-1/2}$ now sit in front, and the $\omega$ integral only has to be
bounded. This is why borderline directions are zoomed by $(\lambda\log t)^{1/2}$
rather than $\lambda^{1/2}$: the $\log t$ that Lemma~\ref{lem:curv} puts into $a_0$
moves into the Jacobian, where nothing can cancel it.

\paragraph{Strips.} Let $\mathcal B_{\mathcal A}$ be the bad directions in
$\mathcal A$: the bad oscillatory ones, and also the borderline ones after an early
switch. Take $\vartheta\in C_0^\infty(\R)$ with $\vartheta=1$ on $[-2,2]$ and
$\vartheta=0$ outside $[-3,3]$. For each $m\in\mathcal B_{\mathcal A}$ write
$1=\vartheta(\omega_m)+(1-\vartheta(\omega_m))$ and expand. This splits
$J=\sum_SJ_S$ over the subsets $S\subset\mathcal B_{\mathcal A}$ of directions placed
in strips, with
\[
J_S=\lambda^{-n/2}\prod_{j\ne k}L_j^{-1/2}\int_{\R^n}e^{i\tilde\Phi}\,\tilde b_S\,d\omega ,
\qquad
\tilde b_S=\tilde b\prod_{m\in S}\vartheta(\omega_m)\prod_{m\in\mathcal B_{\mathcal A}\setminus S}\big(1-\vartheta(\omega_m)\big).
\]
In $J_S$ the \emph{active} directions are $A=\mathcal A\setminus S$. All other
directions, those in $S$ and those outside $\mathcal A$, are \emph{frozen}.

\subsection{Checking the conditions at the zoom scale}

\begin{lemma}[Conditions at the zoom scale]\label{lem:zoomcheck}
	Let $t\ge t_0$ and $S\subset\mathcal B_{\mathcal A}$, and fix the frozen variables.
	View $\tilde\Phi$ and $\tilde b_S$ as functions of $\omega_A\in\R^{|A|}$, and let
	$\tilde K$, $\tilde K_1$, $a_0$, $\tilde M$, $\tilde N$ be the quantities of
	Theorem~\ref{thm:ABZ} for them. Then $\nabla_{\omega_A}\tilde\Phi$ is injective,
	and
	\[
	a_0\ge c,\qquad \tilde M_{|A|+2}\le C,\qquad \tilde N_{|A|+1}\le C,
	\]
	with $c$, $C$ depending only on $d$ and the $\ell_j$.
\end{lemma}

\begin{proof}
	We first locate $\tilde K_1$ in the $\eta$ variables. On $\tilde K$ the amplitude
	does not vanish, so $|\eta|\le2\sqrt d$ and \eqref{eq:support} holds. Passing to
	$\tilde K_1$ moves each active $\eta_j$ by at most $(\lambda L_j)^{-1/2}\le1$. For an
	active bad $m$ this move is
	\[
	(\lambda L_m)^{-1/2}=L_m^{-1/2}\,t^{-\kappa_m/2}|\xi_k|^{1/2}\cdot\frac{t^{\ell_m-\ell_k}}{|\xi_k|}
	\le\sqrt2\,\frac{t^{\ell_m-\ell_k}}{|\xi_k|},
	\]
	because $\kappa_m\ge0$ and $|\xi_k|\le2$. Hence on $\tilde K_1$
	\begin{equation}\label{eq:K1}
	|\eta|\le3\sqrt d,\qquad |\eta_m|\le\frac{4\,t^{\ell_m-\ell_k}}{|\xi_k|}\quad(m\text{ bad}).
	\end{equation}

	\emph{(a) The determinant $a_0$.} By \eqref{eq:K1} and Lemma~\ref{lem:curv} with
	$R=3\sqrt d$, $-\langle\nabla_\eta^2\Phi\,v,v\rangle\ge c\sum_j\Gamma_jv_j^2$ on
	$\tilde K_1$. We claim $\Gamma_j\ge c\,L_j$ for $j\in A$. If $L_j=1$ this is
	Lemma~\ref{lem:curv}(i). If $L_j=\log t$, then $j$ is borderline and the switch is
	early, and \eqref{eq:K1} gives for every bad $m$
	\[
	|\eta_m|^{1/(\ell_k-\ell_m)}\le\Big(\frac4{|\xi_k|}\Big)^{1/(\ell_k-\ell_m)}t^{-1}
	\le t^{-1/2},
	\]
	because $|\xi_k|\ge4t^{-\theta}\ge4t^{-(\ell_k-\ell_m)/2}$. So
	$\sigma_*(\eta)\le t^{-1/2}$, and Lemma~\ref{lem:curv}(ii) gives
	$\Gamma_j\ge\frac12\log t$. Now let $u\in\R^{|A|}$ and put $v_j=L_j^{-1/2}u_j$ for
	$j\in A$ and $v_j=0$ otherwise. By the formula for the Hessian in $\omega$,
	\[
	-\langle\nabla_{\omega_A}^2\tilde\Phi\,u,u\rangle=-\langle\nabla_\eta^2\Phi\,v,v\rangle
	\ge c\sum_{j\in A}\frac{\Gamma_j}{L_j}\,u_j^2\ge c\,|u|^2 .
	\]
	So every eigenvalue of $-\nabla_{\omega_A}^2\tilde\Phi$ is at least $c$, and
	$a_0\ge c^{|A|}$. In the $\eta$ variables this says
	$|\det\nabla_{\eta_A}^2\Phi|\ge c\,(\log t)^{|A\cap\mathcal Z|}$: the log in the
	determinant is exactly the gain $\prod_jL_j^{-1/2}$ in front of $J_S$.

	\emph{(b) Injectivity.} $\tilde\Phi$ is strictly concave in $\omega_A$ on all of
	$\R^{|A|}$, by Section~\ref{sec:hessian} and the linear change of variables. So for
	$\omega\ne\omega'$,
	\[
	\big\langle\nabla\tilde\Phi(\omega)-\nabla\tilde\Phi(\omega'),\omega-\omega'\big\rangle
	=\int_0^1\big\langle\nabla^2\tilde\Phi\big(\omega'+s(\omega-\omega')\big)(\omega-\omega'),\omega-\omega'\big\rangle\,ds<0 ,
	\]
	and $\nabla_{\omega_A}\tilde\Phi$ is injective.

	\emph{(c) The phase derivatives $\tilde M$.} Let $2\le|\alpha|\le|A|+2$ with
	$\alpha$ supported in $A\subset\mathcal O\cup\mathcal Z$, and choose $i,j$ in the
	support of $\alpha$. Then $\partial^\alpha\Phi=-\partial^\alpha G$, and the chain
	rule with Lemma~\ref{lem:weighted} gives
	\[
	|\partial_\omega^\alpha\tilde\Phi|
	\le C\Big(\frac{\bar\Lambda_i\bar\Lambda_j}{L_iL_j}\Big)^{1/2}
	\prod_{m\ne k}\big((\lambda L_m)^{1/2}v_m\big)^{-\alpha'_m}.
	\]
	The first factor is bounded, since $\bar\Lambda_m\le C$ when $L_m=1$ and
	$\bar\Lambda_m=\log t=L_m$ for active borderline $m$. In the product, a good $m$
	has $(\lambda L_m)^{1/2}v_m\ge\lambda^{1/2}\ge1$. An active bad $m$ lies in
	$\mathcal B_{\mathcal A}\setminus S$ and has $(\lambda L_m)^{1/2}v_m=|\omega_m|$.
	Since $1-\vartheta(\omega_m)$ vanishes for $|\omega_m|\le2$ and $\tilde K_1$ extends
	the support by sup distance $1$, we have $|\omega_m|\ge1$ on $\tilde K_1$. So every
	factor is at most $1$, and $\tilde M_{|A|+2}\le C$.

	\emph{(d) The amplitude derivatives $\tilde N$.} One $\omega_j$ derivative, $j\in
	A$, is $(\lambda L_j)^{-1/2}\partial_{\eta_j}$. Falling on $(1+|\eta|^2)^{-1/2}$ or
	on $\rho_k$, it costs at most $C(\lambda L_j)^{-1/2}\le C$. Falling on
	$\tilde\chi$, it produces the factor $t^{\ell_k-\ell_j}\xi_k$, so it costs
	\[
	(\lambda L_j)^{-1/2}\,t^{\ell_k-\ell_j}|\xi_k|=L_j^{-1/2}\,t^{-\kappa_j/2}\,|\xi_k|^{1/2}
	\le\sqrt2 ,
	\]
	because $\kappa_j\ge0$, $L_j\ge1$ and $|\xi_k|\le2$. This is where $\kappa_j\ge0$
	is used. Derivatives of the strip cutoffs are $O(1)$. So
	$\tilde N_{|A|+1}\le C$.
\end{proof}

The check in one table:

\begin{center}
\begin{tabular}{lll}
	& unit scale in $\eta$ & zoom scale in $\omega$\\\hline
	$a_0$ & $\ge c\,(\log t)^{|A\cap\mathcal Z|}$, offset by $M$ & $\ge c$, gain moved to the Jacobian\\
	$M$ & grows with $t$ in bad directions & $\le C$ off the strips\\
	$N$ & up to $t^{d(\ell_k-\ell_{\min})}$ & $\le C$ in $\mathcal O\cup\mathcal Z$\\
\end{tabular}
\end{center}

\section{Proof of the main estimate}\label{sec:proof}

\subsection{Statement}

Recall from Section~\ref{sec:spt} that $b(\xi_k,\cdot)$ is smooth and supported in
$|\eta|\le2\sqrt d$, and that $|\xi_k|\le2$ and $\lambda\ge c\,t^{1+\ell_{\min}}$
whenever $b(\xi_k,\cdot)$ is not identically zero.

\begin{theorem}[Main estimate]\label{thm:main}
	There are $t_0$ and $C$, depending only on $d$ and $\ell_1,\dots,\ell_d$, such
	that for $t\ge t_0$, all $x\in\R^d$ and all $\xi_k\ne0$,
	\[
	\Big|\int_{\R^{d-1}}e^{i\lambda\Phi(\eta)}\,b(\xi_k,\eta)\,d\eta\Big|
	\le C\,\lambda^{-\frac{d-1}{2}}\prod_{j\in\mathcal F}t^{\kappa_j/2}\,|\xi_k|^{-1/2}
	\prod_{j\in\mathcal Z}E(\xi_k),
	\qquad\lambda=t^{\ell_k+1}|\xi_k| ,
	\]
	where $E(\xi_k)=(\log t)^{-1/2}$ if $|\xi_k|\ge4t^{-\theta}$ and
	$E(\xi_k)=|\xi_k|^{-1/2}$ otherwise.
\end{theorem}

Each factor in the product over $\mathcal F$ is at most $|\xi_k|^{-1/2}$, since
$\kappa_j<0$ there. When every direction is oscillatory, the estimate is the full
stationary phase rate $\lambda^{-(d-1)/2}$. Each borderline direction gives an extra
factor $(\log t)^{-1/2}$ after an early switch. After a late switch it is treated
like a flat direction with $\kappa_j=0$.

\begin{remark}
	The factors $|\xi_k|^{-1/2}$ have to be handled in the outer $\xi_k$ integral of
	$I_k$. So does the late-switch region $|\xi_k|<4t^{-\theta}$, which has length at
	most $8t^{-\theta}$.
\end{remark}

\subsection{Proof}

Write $J$ for the integral in Theorem~\ref{thm:main}, with $t$, $x$, $\xi_k$ fixed. If
$b(\xi_k,\cdot)\equiv0$ there is nothing to prove, so we assume $|\xi_k|\le2$ and
$\lambda\ge c\,t^{1+\ell_{\min}}$, and we take $t_0\ge e$ so large that
$\lambda\ge1$.

\paragraph{Step 1: zoom and strips.} By Section~\ref{sec:zoom},
$J=\sum_SJ_S$ over the $2^{|\mathcal B_{\mathcal A}|}$ subsets
$S\subset\mathcal B_{\mathcal A}$, with the decay
$\lambda^{-n/2}\prod_jL_j^{-1/2}$ in front of each $J_S$.

\paragraph{Step 2: stationary phase in the active variables.} Fix $S$ and the frozen
variables. By Lemma~\ref{lem:zoomcheck}, Theorem~\ref{thm:ABZ} applies to the
$\omega_A$ integral with $\lambda=1$ and gives
\[
\Big|\int_{\R^{|A|}}e^{i\tilde\Phi}\,\tilde b_S\,d\omega_A\Big|\le C ,
\]
uniformly in $x$ and in the frozen variables.

\paragraph{Step 3: trivial integration in the frozen variables.}
\begin{itemize}
	\item $m\in S$. The cutoff $\vartheta(\omega_m)$ is supported in $|\omega_m|\le3$,
	so this costs a factor $6$.
	\item $j\notin\mathcal A$, a flat direction or a borderline one after a late
	switch. Then $L_j=1$, and \eqref{eq:support} gives $|\omega_j|\le W_j$ with
	\[
	W_j=\lambda^{1/2}\cdot\frac{2t^{\ell_j-\ell_k}}{|\xi_k|}
	=2\,t^{\kappa_j/2}\,|\xi_k|^{-1/2}\le2\,|\xi_k|^{-1/2},
	\]
	so this costs a factor $2W_j$. Here $\kappa_j\le0$ is used: the support in
	$\omega_j$ is no wider than a constant times $|\xi_k|^{-1/2}$.
\end{itemize}

\paragraph{Step 4: conclusion.} Each piece satisfies
\[
|J_S|\le C\,\lambda^{-\frac{d-1}{2}}\prod_{j\ne k}L_j^{-1/2}\;6^{|S|}\prod_{j\notin\mathcal A}2W_j .
\]
After an early switch, $\prod_jL_j^{-1/2}=(\log t)^{-|\mathcal Z|/2}$ and the
directions outside $\mathcal A$ are the flat ones. After a late switch, all $L_j=1$
and the directions outside $\mathcal A$ are the flat and the borderline ones, with
$t^{\kappa_j/2}=1$ for the latter. Summing over $S$ proves Theorem~\ref{thm:main}.
The bound is uniform in $x$ because $x$ enters $\Phi$ only through its linear part,
which none of $a_0$, $\tilde M$, $\tilde N$ sees.

\paragraph{Summary.} After zooming to $\omega_j=(\lambda L_j)^{1/2}\eta_j$, the
condition $\kappa_j\ge0$ is exactly what keeps the amplitude smooth at scale $1$, and
Lemma~\ref{lem:curv} keeps the Hessian uniformly nondegenerate. Oscillatory directions
gain $\lambda^{-1/2}$. Borderline directions gain $(\lambda\log t)^{-1/2}$ after an
early switch, because every dyadic time scale after the switch adds the same amount
of curvature. Flat directions have support thinner than the stationary phase scale
and are integrated trivially. In bad directions the weighted bounds degenerate on a
strip of unit width in $\omega$, and integrating over the strip costs the same factor
that oscillation would gain.

\end{document}
